\paragraph{Pauli-Z (Phase Flip) Gate}\label{sec:pauli-z-phase-flip-gate}

The \definitionWithSpecificIndex{Pauli-Z Gate}{Single Qubit Gates: Pauli-Z Gate}{}, also called the \textbf{Phase Flip Gate}, is one of the core \textbf{Pauli operators} used in quantum computing. It acts \textbf{only on the phase} of the qubit, \textbf{leaving the \hl{$\ket{0}$ amplitude unchanged}} but \textbf{\hl{flipping} the sign of the \hl{$\ket{1}$ amplitude}}.

\highspace
\begin{flushleft}
    \textcolor{Green3}{\faIcon{square-root-alt} \textbf{Matrix Representation}}
\end{flushleft}
\begin{equation}
    Z = \sigma_{Z} = 
    \begin{pmatrix}
        1 & 0 \\
        0 & -1
    \end{pmatrix}
\end{equation}

\highspace
\begin{flushleft}
    \textcolor{Green3}{\faIcon{book} \textbf{Action on a General Qubit}}
\end{flushleft}
For a general state:
\begin{equation*}
    \ket{\psi} = a\ket{0} + b\ket{1}
\end{equation*}
Applying $Z$:
\begin{equation*}
    Z\ket{\psi} = aZ\ket{0} + bZ\ket{1} = a\ket{0} - b\ket{1}
\end{equation*}
As a result, $Z$ \textbf{flips the sign of the $\ket{1}$ amplitude}, this is why it's called a \textbf{phase flip}. Finally, it's trivial to show that:
\begin{equation*}
    Z\ket{0} = \ket{0} \hspace{1em} \text{(unchanged)}, \hspace{2em} Z\ket{1} = -\ket{1} \hspace{1em} \text{(sign flip)}
\end{equation*}
So the \textbf{magnitudes} of the amplitudes remain unchanged and the \textbf{probabilities of measurement outcomes} in the standard basis are unaffected, but the \textbf{relative phase} between $\ket{0}$ and $\ket{1}$ is altered (since $b$ changes to $-b$).

\begin{examplebox}[: Effect on Superposition States]
    The Hadamard (or X-basis) states are:
    \begin{equation*}
        \ket{+} = \dfrac{\ket{0} + \ket{1}}{\sqrt{2}}, \hspace{2em} \ket{-} = \dfrac{\ket{0} - \ket{1}}{\sqrt{2}}
    \end{equation*}
    These two states differ \textbf{only by their relative phase}:
    \begin{itemize}
        \item $\ket{+}$: relative phase $0$.
        \item $\ket{-}$: relative phase $\pi$.
    \end{itemize}
    Applying $Z$ on $\ket{+}$:
    \begin{align*}
        Z\ket{+} =& Z\left( \dfrac{\ket{0} + \ket{1}}{\sqrt{2}} \right) \\
        =& \dfrac{Z\ket{0} + Z\ket{1}}{\sqrt{2}} \\
        =& \dfrac{\ket{0} - \ket{1}}{\sqrt{2}} \\
        =& \ket{-}
    \end{align*}
    Applying $Z$ on $\ket{-}$:
    \begin{align*}
        Z\ket{-} =& Z\left( \dfrac{\ket{0} - \ket{1}}{\sqrt{2}} \right) \\
        =& \dfrac{Z\ket{0} - Z\ket{1}}{\sqrt{2}} \\
        =& \dfrac{\ket{0} + \ket{1}}{\sqrt{2}} \\
        =& \ket{+}
    \end{align*}
    The Pauli-Z gate \textbf{does not change the probabilities} of measuring $\ket{0}$ or $\ket{1}$ in the computational basis. However, it \textbf{changes the relative phase}, transforming $\ket{+}$ into $\ket{-}$ and vice versa.
\end{examplebox}

\begin{flushleft}
    \textcolor{Green3}{\faIcon{question-circle} \textbf{Geometric Interpretation: Rotation around the $z$-axis}}
\end{flushleft}
On the Bloch sphere, the Pauli-Z gate performs a \textbf{$\pi$ radians ($180\degree$) rotation around the $z$-axis}. In spherical coordinates:
\begin{equation*}
    \ket{\psi} = \cos\left(\dfrac{\theta}{2}\right)\ket{0} + e^{i \phi} \sin\left(\dfrac{\theta}{2}\right)\ket{1}
\end{equation*}
After Z rotation:
\begin{itemize}
    \item The angle $\phi \rightarrow \phi + \pi$
    \item This phase shift \textbf{changes the sign of the complex amplitude for} $\ket{1}$
\end{itemize}
Since \textbf{measurement probabilities depend on} $\left|a\right|^{2}$ and $\left|b\right|^{2}$, which \hl{remain unchanged}, the Z gate \textbf{does not affect measurement outcomes} in the standard basis. It only affects \textbf{interference and future gate interactions}.

\begin{figure}[!htp]
    \centering
    \tdplotsetmaincoords{70}{120}
    \begin{tikzpicture}[tdplot_main_coords, scale=3.4, >=Latex]

        \def\r{1}

        % Axes
        \draw[->, thick] (0,0,0) -- (1.25,0,0) node[anchor=north east] {$x$};
        \draw[->, thick] (0,0,0) -- (0,1.25,0) node[anchor=west] {$y$};
        \draw[->, very thick, Green3] (0,0,-1.2) -- (0,0,1.35) node[anchor=south] {$z$};

        % Main sphere circle
        \tdplotdrawarc[dashed]{(0,0,0)}{\r}{0}{180}{}{}
        \tdplotdrawarc[thick]{(0,0,0)}{\r}{180}{360}{}{}

        % Vertical great circle
        \tdplotsetrotatedcoords{90}{0}{90}
        \tdplotdrawarc[tdplot_rotated_coords,dashed]{(0,0,0)}{\r}{0}{180}{}{}
        \tdplotdrawarc[tdplot_rotated_coords,thick]{(0,0,0)}{\r}{180}{360}{}{}
        \tdplotsetrotatedcoords{0}{0}{0}

        % Eigenstates
        \fill[Green3] (0,0,1) circle (0.018);
        \node[Green3, anchor=south east] at (0,0,1.08) {$\ket{0}$};

        \fill[Green3] (0,0,-1) circle (0.018);
        \node[Green3, anchor=north east] at (0,0,-1.08) {$\ket{1}$};

        % Rotation around z-axis
        \tdplotdrawarc[->, very thick, magenta]
            {(0,0,0)}{0.72}{20}{200}
            {anchor=south}{$\pi$ rotation around $z$}
    \end{tikzpicture}
    \caption{Eigenstates of the Pauli-$Z$ gate. Since $\ket{0}$ and $\ket{1}$ lie on the $z$-axis, they are not moved by a rotation around that axis. The state $\ket{1}$ only acquires a global phase $-1$.}
    \label{fig:pauli-z-eigenstates}
\end{figure}

\highspace
\begin{flushleft}
    \textcolor{Green3}{\faIcon{book} \textbf{Interpretation and Use Cases}}
\end{flushleft}
\begin{itemize}
    \item \important{Phase Shift}: $Z$ is a special case of a phase shift gate, shifting the phase by $\pi$.
    \item \important{No effect on $\ket{0}$}: Useful for controlled operations, where only $\ket{1}$ paths acquire a phase.
\end{itemize}

\highspace
\begin{flushleft}
    \textcolor{Green3}{\faIcon{list-ul} \textbf{Key Properties}}
\end{flushleft}
\begin{enumerate}
    \item \textbf{Unitary}: $Z^{\dagger} = Z^{-1} = Z$. Ensuring normalization is preserved.
    \item \textbf{Hermitian}: $Z^{\dagger} = Z$. Meaning it is equal to its own conjugate transpose.
    \item \textbf{Self-Inverse}: $Z^{-1} = Z \Rightarrow Z^{2} = I$. Applying the Pauli-Z gate twice returns the original state.
    \item \textbf{Eigenvalues and Eigenstates}. The Pauli-Z operator has \emph{eigenvalues} $+1$ and $-1$, with corresponding \emph{eigenstates} $\ket{0}$ and $\ket{1}$:
    \begin{equation*}
        Z\ket{0} = +1 \cdot \ket{0}, \hspace{2em} Z\ket{1} = -1 \cdot \ket{1}
    \end{equation*}
    This means that $\ket{0}$ and $\ket{1}$ are \textbf{unchanged in direction} by the action of $Z$ (up to a possible phase factor). In particular:
    \begin{itemize}
        \item $\ket{0}$ is completely unchanged.
        \item $\ket{1}$ acquires only a phase factor $-1 = e^{i\pi}$.
    \end{itemize}
    Therefore, these states are \textbf{invariant under the action of $Z$} (up to global phase), and represent the \textbf{natural states of the operator}.
\end{enumerate}

\highspace
\textcolor{Green3}{\faIcon[regular]{lightbulb} \textbf{Geometric Insight.}} On the Bloch sphere, the Pauli-Z gate corresponds to a rotation around the $z$-axis. The eigenstates $\ket{0}$ and $\ket{1}$ lie exactly on this axis (north and south poles). As a consequence, they are \textbf{not affected by the rotation}, which explains why they are eigenstates of $Z$.

\highspace
In contrast, any state that is not aligned with the $z$-axis (for example $\ket{+}$, \autoref{fig:bloch-sphere}) is rotated by the action of $Z$. This highlights the difference between:
\begin{itemize}
    \item \textbf{Eigenstates}: unchanged (up to phase)
    \item \textbf{Superposition states}: geometrically rotated
\end{itemize}